% ---------------- 第3章 章末 ----------------
\zhangmo{本章小结}
\subsection*{一、知识总结}
\begin{center}
\small
\begin{tabularx}{\linewidth}{p{2.4cm}X}
\toprule
平衡状态 & 逆、等、动、定、变；判断：正逆相等（一正一逆、成比例）或变量不变 \\
平衡常数 & 只与温度有关；固体纯液体不写；方程式相加 $K$ 相乘；$Q$ 与 $K$ 比较定方向；$K_p$ 用分压 \\
平衡移动 & 浓度、压强（看气体分子数）、温度（唯一改变 $K$）、催化剂（不移动）；勒夏特列原理“减弱不抵消” \\
图像 & 四看；先拐先平数值大；定一议二；曲线外的点看趋向；最高点前未平衡、后平衡 \\
方向 & $\Delta G=\Delta H-T\Delta S<0$ 自发；四种组合；转折温度 $\Delta H/\Delta S$ \\
调控 & 合成氨 10～30 MPa、400～500 ℃、铁触媒、液化分离与循环；兼顾速率、平衡、成本 \\
多重平衡 & 每个反应设变量；原子守恒；同一体系共用浓度；选择性$=$生成目标产物所耗反应物/反应物总转化量 \\
\bottomrule
\end{tabularx}
\end{center}

\subsection*{二、方法提升}
\begin{enumerate}
\item \textbf{三段式是平衡计算的“万能钥匙”}：起始、转化、平衡三行，转化量与计量数成比例；有 $K_p$ 时加一行“分压”。
\item \textbf{判断移动方向先用 $Q$ 与 $K$}：浓度、压强变化看 $Q$ 怎么变，温度变化看 $K$ 怎么变。
\item \textbf{“假设不移动再比较”}：处理“浓度变为原来的几倍”类问题，先算假设不移动时的值，再与实际值比较。
\item \textbf{多重平衡温度曲线}：先判断各反应的吸放热，再分温度段讨论谁占主导。
\end{enumerate}

\subsection*{三、易错点归纳}
\begin{enumerate}
\item 把“浓度相等”“浓度比等于计量数比”当成平衡标志；忘记检查一个量是否为“变量”。
\item $K$ 表达式中写入固体、纯液体；用起始浓度或某时刻浓度代替平衡浓度计算 $K$。
\item 误认为增大压强平衡一定移动；恒容充惰性气体判为“加压，平衡移动”。
\item 误认为平衡正移，反应物转化率一定增大（增大某反应物浓度时其自身转化率减小）。
\item 误认为“自发反应一定能迅速发生”；计算 $T=\Delta H/\Delta S$ 时忘统一 kJ 与 J。
\item 选择合成氨温度时答“温度高有利于平衡正移”（错，低温才有利于平衡）。
\end{enumerate}

\zhangmo{分层练习}
\zu{A组\quad 基础巩固}
\begin{xiti}
\item 一定温度下的恒容容器中发生 \ce{2A(g) + B(s) <=> C(g) + D(g)}，下列描述能表明反应已达平衡的是\xz\\
①混合气体的压强不变\quad ②混合气体的密度不变\quad ③C(g) 的浓度不变\quad ④A、C、D 的浓度之比为 2:1:1\quad ⑤单位时间内生成 $n$ mol C，同时生成 $n$ mol D\quad ⑥单位时间内生成 $n$ mol D，同时生成 $2n$ mol A
\xx{①②⑤}{②③⑤}{②③⑥}{①③⑥}

\item 下列关于化学平衡常数的说法正确的是\xz
\xxa{化学平衡常数表示反应进行的程度，与温度无关}{两种物质反应，不管怎样书写化学方程式，平衡常数均不变}{化学平衡常数等于某一时刻生成物浓度幂之积与反应物浓度幂之积的比值}{温度一定时，对于给定的反应，正、逆反应的平衡常数互为倒数}

\item 某温度下 \ce{N2(g) + 3H2(g) <=> 2NH3(g)} 的平衡常数为 $K$。同温下，\ce{1/2N2(g) + 3/2H2(g) <=> NH3(g)} 的平衡常数为\tk，\ce{2NH3(g) <=> N2(g) + 3H2(g)} 的平衡常数为\tk。

\item 下列事实\textbf{不能}用勒夏特列原理解释的是\xz
\xxa{光照新制的氯水时，溶液的 pH 逐渐减小}{加催化剂，使 \ce{N2} 和 \ce{H2} 在一定条件下转化为 \ce{NH3}}{可用浓氨水和氢氧化钠固体快速制氨}{增大压强，有利于 \ce{SO2} 与 \ce{O2} 反应生成 \ce{SO3}}

\item 对于 \ce{2SO2(g) + O2(g) <=> 2SO3(g)} $\Delta H<0$，在恒温恒容条件下达到平衡后，下列措施能提高 \ce{SO2} 平衡转化率的是\xz\\
①升高温度\quad ②再充入 \ce{O2}\quad ③充入 He\quad ④再按原比例充入 \ce{SO2} 和 \ce{O2}\quad ⑤使用催化剂\quad ⑥移走部分 \ce{SO3}
\xx{①②⑥}{②④⑥}{②③④}{④⑤⑥}

\item 下列反应在任何温度下均能自发进行的是\xz
\xxa{\ce{2N2(g) + O2(g) = 2N2O(g)} \quad $\Delta H=+163$ kJ·mol$^{-1}$}{\ce{Ag(s) + 1/2Cl2(g) = AgCl(s)} \quad $\Delta H=-127$ kJ·mol$^{-1}$}{\ce{HgO(s) = Hg(l) + 1/2O2(g)} \quad $\Delta H=+91$ kJ·mol$^{-1}$}{\ce{H2O2(l) = 1/2O2(g) + H2O(l)} \quad $\Delta H=-98$ kJ·mol$^{-1}$}

\item 在 \ce{2NO2(g) <=> N2O4(g)} $\Delta H<0$ 的平衡体系中，下列操作能使混合气体颜色变深、且平衡逆向移动的是\xz
\xxa{恒温时压缩容器体积}{恒容时升高温度}{恒容时充入 Ar}{恒压时充入 Ar}
\end{xiti}

\zu{B组\quad 高考典型}
\begin{xiti}
\item 某温度下，在 2 L 密闭容器中充入气体 A、B，其浓度分别为 2 mol·L$^{-1}$、1 mol·L$^{-1}$，发生 \ce{3A(g) + 2B(g) <=> 4C(\text{?}) + 2D(\text{?})}。达到平衡时生成 1.6 mol C，且反应前后的压强之比为 5:4（相同温度）。(1) 判断 C、D 的状态，写出平衡常数表达式；(2) 求 B 的平衡转化率和该温度下的 $K$（保留两位有效数字）；(3) 增加 C 的量，A、B 的转化率如何变化？

\item \laiyuan{原创}\ 对于 \ce{$m$A(g) + $n$B(g) <=> $p$C(g)}，在不同条件下 C 的体积分数随时间变化如下：曲线 a（$T_2$、$p_2$）最先出现拐点，平台高度 0.30；曲线 b（$T_1$、$p_2$）其次出现拐点，平台高度 0.50；曲线 c（$T_1$、$p_1$）最后出现拐点，平台高度 0.35。(1) 比较 $T_1$、$T_2$ 和 $p_1$、$p_2$ 的大小；(2) 判断 $\Delta H$ 的正负及 $m+n$ 与 $p$ 的关系；(3) 在恒压下，\ce{2SO2(g) + O2(g) <=> 2SO3(g)} 的 \ce{SO2} 平衡转化率随温度升高而降低，投料比 $m=n(\ce{SO2})/n(\ce{O2})$ 分别为 $m_1$、$m_2$、$m_3$ 的三条曲线自上而下排列，比较 $m_1$、$m_2$、$m_3$ 的大小。

\item 硫酸工业中 \ce{2SO2(g) + O2(g) <=> 2SO3(g)} $\Delta H=-196.6$ kJ·mol$^{-1}$。不同条件下 \ce{SO2} 的平衡转化率：450 ℃ 时 0.1 MPa 为 97.5\%，10 MPa 为 99.7\%；550 ℃ 时 0.1 MPa 为 85.6\%，10 MPa 为 98.3\%。(1) 实际生产为何选择常压？(2) 为何选择 400～500 ℃ 而不是更低温度？(3) 为何通入过量空气？(4) 尾气为何要处理？

\item \laiyuan{2024年广东卷第17题(4)，来源：仓库所附真题文档}\ 兴趣小组探究反应 \ce{CaSO4(s) + CO3^2-(aq) <=> CaCO3(s) + SO4^2-(aq)}（反应Ⅰ）平衡的建立：将 1.00 g 硫酸钙（$M=136$ g·mol$^{-1}$）加入 100.0 mL 0.100 mol·L$^{-1}$ \ce{Na2CO3} 溶液中，在 25 ℃ 和搅拌条件下，用 pH 计测得体系 pH 随时间先减小，$t_2$ min 后不变。甲同学认为反应Ⅰ已达平衡；乙同学认为证据不足，提出：假设1 硫酸钙固体已完全消耗；假设2 硫酸钙有剩余，但被碳酸钙沉淀包裹。(1) 乙设计实验：将反应混合物固液分离，取少量沉淀置于试管中，滴加\tk，现象为\tk，沉淀完全溶解；继续向试管中滴加\tk，无白色沉淀生成。由此得出假设1成立。(2) 从数据角度说明假设1为何成立。(3) 写出优化的实验方案，并给出反应Ⅰ平衡已建立的判断依据。

\item \laiyuan{原创}\ 在恒定总压 20 MPa、一定温度下，向反应器中充入 1 mol \ce{N2} 和 3 mol \ce{H2}，发生 \ce{N2(g) + 3H2(g) <=> 2NH3(g)}，平衡时 \ce{NH3} 的物质的量分数为 25\%。(1) 求 \ce{H2} 的平衡转化率。(2) 求该温度下的 $K_p$（保留三位有效数字，单位 MPa$^{-2}$）。(3) 若改为在恒温恒容条件下进行（起始压强仍为 20 MPa），\ce{NH3} 的平衡物质的量分数将\tk（填“大于”“小于”或“等于”）25\%。

\item \laiyuan{原创}\ \ce{CO2(g) + 3H2(g) <=> CH3OH(g) + H2O(g)} $\Delta H<0$。在压强一定时，按不同投料比 $m=\dfrac{n(\ce{H2})}{n(\ce{CO2})}$（$m_1$、$m_2$、$m_3$）进行实验，\ce{CO2} 的平衡转化率随温度变化的三条曲线均随温度升高而下降，且同一温度下 $m_1$ 对应的转化率最高、$m_3$ 最低。(1) 比较 $m_1$、$m_2$、$m_3$ 的大小。(2) 解释曲线随温度升高而下降的原因。(3) 若同一投料比下，曲线上 a 点（$T_a$）与 b 点（$T_b$，$T_b>T_a$）的平衡常数分别为 $K_a$、$K_b$，比较大小。
\end{xiti}

\zu{C组\quad 冲刺高分}
\begin{xiti}
\item \laiyuan{原创}\ 乙苯脱氢制苯乙烯：\ce{C6H5CH2CH3(g) <=> C6H5CH=CH2(g) + H2(g)} $\Delta H>0$。某温度、总压恒为 100 kPa 时，以纯乙苯为原料，乙苯的平衡转化率为 40\%。(1) 求该温度下的 $K_p$。(2) 保持温度和总压不变，改为通入 $n(\ce{H2O}):n(\text{乙苯})=9:1$ 的混合气（水蒸气不参与反应），求乙苯的平衡转化率（保留两位有效数字）。(3) 解释工业上掺入大量水蒸气的原因（至少答两点）。

\item \laiyuan{2025年广东卷第19题(4)，来源：仓库所附真题文档}\ 钛基催化剂可催化储氢物质肼（\ce{N2H4}）的分解：(e)\ce{N2H4 = N2 + 2H2}；(f)\ce{3N2H4 = N2 + 4NH3}。以肼的水溶液为原料（含 \ce{N2H4} 的物质的量为 $n_0$）进行实验，得到 $n_1/n_0$、$n_2/n_0$ 随时间 $t$ 的变化曲线，其中 $n_1$ 为 \ce{H2} 与 \ce{N2} 的物质的量之和，$n_2$ 为剩余 \ce{N2H4} 的物质的量；$t_0$ min 时 $n_1/n_0=y_1$，$n_2/n_0=y_2$。设 $n_\text{e}$ 为 $0\sim t$ 时间段内反应(e)消耗 \ce{N2H4} 的物质的量，该时间段内催化剂的选择性用 $\dfrac{n_\text{e}}{n_0-n_2}\times100\%$ 表示。①$0\sim t_0$ min 内，\ce{N2H4} 的转化率为\tk（用含 $y_2$ 的代数式表示）。②$0\sim t_0$ min 内，催化剂的选择性为\tk（用含 $y_1$ 与 $y_2$ 的代数式表示，写出推导过程）。

\item \laiyuan{2026年广东卷第19题(5)，来源：仓库所附真题文档（网络收集版）}\ 某低共熔溶剂 D(l) 可用于吸收 HCl(g)，吸收过程可表示为 \ce{HCl(g) <=> HCl(sol)}，\ce{$x$HCl(sol) + D(l) <=> D*$x$HCl(sol)}。在一定 $T$ 和 HCl(g) 压强 $p$ 下，吸收至饱和时，每摩尔 D 吸收 HCl(g) 的量为 $\alpha+\beta$，其中 $\alpha=\dfrac{n(\ce{HCl})}{n_0(\mathrm{D})}$，$\beta=\dfrac{x\cdot n(\mathrm{DH})}{n_0(\mathrm{D})}$（$x$ 为仅与温度有关的系数；$n_0(\mathrm{D})$ 为 D(l) 的初始物质的量；$n(\ce{HCl})$ 为平衡时 HCl(sol) 的物质的量；$n(\mathrm{DH})$ 为平衡时 \ce{D*$x$HCl(sol)} 的物质的量）。由图读得：303 K、$p=100$ kPa 时 $\alpha=0.15$、$\beta=0.86$；350 K 时 $\alpha$ 始终为 0.00，$p=40$ kPa 时 $\beta=0.19$，$p=100$ kPa 时 $\beta=0.25$。
①1 mol D 在 303 K、$p=100$ kPa 下吸收 HCl(g) 至饱和，再升温至 350 K、降压至 $p=40$ kPa，达平衡后可释放\tk mol HCl(g)。
②形成 \ce{D*$x$HCl(sol)} 的总过程为 \ce{$x$HCl(g) + D(l) <=> D*$x$HCl(sol)}，其平衡常数 $K=\dfrac{n(\mathrm{DH})}{n(\mathrm{D})\cdot(p/100\,\mathrm{kPa})^x}$。已知 350 K 时 $K=1.0$，则在 350 K、$p=64$ kPa 下，1 mol D 最多可吸收\tk mol HCl(g)（写出计算过程，结果保留两位有效数字）。

\item \laiyuan{原创}\ 一定温度下，在恒容密闭容器中充入 \ce{NO2}，起始压强为 100 kPa，发生 \ce{2NO2(g) <=> N2O4(g)}，该反应可视为基元反应：$\vz=\kz p^2(\ce{NO2})$，$\vn=\kn p(\ce{N2O4})$。达到平衡时总压为 70 kPa。(1) 求平衡时 \ce{NO2}、\ce{N2O4} 的分压及 $K_p$。(2) 用 $\kz$、$\kn$ 表示 $K_p$。(3) 当 \ce{NO2} 的分压为 60 kPa 时，$\dfrac{\vz}{\vn}=$\tk。(4) 升高温度，平衡体系中 \ce{NO2} 的体积分数增大，判断该反应的 $\Delta H$ 正负。

\item \laiyuan{改编自选必1课件“演示文稿1”中的练习，数据为本书设定}\ 羰基法提纯镍：\ce{Ni(s) + 4CO(g) <=> Ni(CO)4(g)}。在不同总压下，平衡时 \ce{Ni(CO)4} 的体积分数 $x$ 随温度升高而减小，且同一温度下总压越大 $x$ 越大。(1) 判断该反应的 $\Delta H$\tk0，从热力学角度考虑，\tk{}有利于 \ce{Ni(CO)4} 的生成（写出两点）。(2) 设起始只充入 CO（Ni 足量），若总压为 1.0 MPa、100 ℃ 时 $x=0.90$，求 CO 的平衡转化率（保留三位有效数字）和该温度下的 $K_p$（单位 MPa$^{-3}$）。(3) 工业上先在约 50 ℃ 使粗镍与 CO 反应生成 \ce{Ni(CO)4}，再将气体送到约 230 ℃ 的分解器得到纯镍。结合 (1) 说明其原理。
\end{xiti}

\zhangmo{练习解析}
\begin{jiexilist}
\item \ans{C}B 为固体，气体分子数 $2\to2$ 不变，恒容时压强始终不变，①不能作标志；气体质量随反应变化，密度是变量，②能；③浓度不变，能；④浓度比不能说明不变；⑤都表示正反应；⑥生成 D 为正、生成 A 为逆，$n:2n=1:2$ 与计量数比相符，能。选②③⑥。
\item \ans{D}$K$ 与温度有关，A 错；方程式写法不同 $K$ 不同，B 错；$K$ 用平衡浓度，C 错；正逆反应 $K$ 互为倒数，D 对。
\item \ans{$K^{1/2}$（即 $\sqrt K$）；$1/K$}计量数变为 $\frac12$，$K$ 开平方；逆反应取倒数。
\item \ans{B}催化剂不使平衡移动，不能用勒夏特列原理解释。A：\ce{HClO} 见光分解，\ce{Cl2 + H2O <=> H+ + Cl- + HClO} 正移，$c(\ce{H+})$ 增大；C：NaOH 溶解放热且 $c(\ce{OH-})$ 增大，使 \ce{NH3 + H2O <=> NH3*H2O <=> NH4+ + OH-} 逆移；D：加压平衡向气体分子数减小的方向移动。
\item \ans{B}①升温平衡逆移，转化率降低；②增大 \ce{O2} 浓度，\ce{SO2} 转化率提高；③恒容充 He，平衡不移动；④按比例加料相当于加压，平衡正移，转化率提高；⑤催化剂不影响平衡；⑥移走生成物，平衡正移。选②④⑥。
\item \ans{D}任何温度自发需 $\Delta H<0$ 且 $\Delta S>0$。A：吸热、气体分子数减少（熵减），任何温度不自发；B：放热、熵减，低温自发；C：吸热、熵增，高温自发；D：放热、生成气体熵增，任何温度自发。
\item \ans{B}A：压缩体积，\ce{NO2} 浓度瞬间增大，颜色变深，但平衡正向（气体分子数减小的方向）移动；B：升温，平衡向吸热的逆反应方向移动，\ce{NO2} 增多，颜色变深，正确；C：恒容充 Ar，各物质浓度不变，平衡不移动，颜色不变；D：恒压充 Ar，体积增大，相当于减压，平衡逆向移动，但 \ce{NO2} 浓度减小，颜色变浅。
\item (1) 起始 A 4 mol、B 2 mol（共 6 mol 气体）；转化 A 1.2 mol、B 0.8 mol、C 1.6 mol、D 0.8 mol。平衡气体总量 $6\times\frac45=4.8$ mol。若 C、D 均为气体，平衡时 $6-2.0+2.4=6.4$ mol，不符；若只有 D 为气体，$6-2.0+0.8=4.8$ mol，符合。故 C 为非气态（固体或液体），D 为气体，$K=\dfrac{c^2(\mathrm{D})}{c^3(\mathrm{A})\cdot c^2(\mathrm{B})}$。(2) $\alpha(\mathrm{B})=0.8/2=40\%$；平衡浓度 $c(\mathrm{A})=1.4$、$c(\mathrm{B})=0.6$、$c(\mathrm{D})=0.4$ mol·L$^{-1}$，$K=\dfrac{0.4^2}{1.4^3\times0.6^2}\approx0.16$。(3) C 非气态，增加 C 的量平衡不移动，A、B 转化率不变。
\item (1) 曲线 a、b 压强相同，a 先拐，$T_2>T_1$；曲线 b、c 温度相同，b 先拐，$p_2>p_1$。(2) 同为 $p_2$，高温（a）C\% 低，升温平衡逆移，$\Delta H<0$；同为 $T_1$，高压（b）C\% 高，加压平衡正移，$m+n>p$。(3) 投料比 $m$ 越小，\ce{O2} 相对越多，\ce{SO2} 转化率越高，最上方曲线 $m$ 最小：$m_1<m_2<m_3$。
\item (1) 常压下 \ce{SO2} 平衡转化率已很高（450 ℃ 时 97.5\%），加压提高有限，而对设备要求和能耗大幅增加，不经济。(2) 温度过低反应速率太慢，且催化剂（\ce{V2O5}）在 400～500 ℃ 活性最高；温度过高平衡转化率下降。(3) 增大 \ce{O2} 浓度，使平衡正移，提高 \ce{SO2} 的转化率（空气廉价易得）。(4) \ce{SO2} 有毒、会形成酸雨污染环境，回收还可提高原料利用率。
\item (1) 过量稀盐酸；有气体产生；\ce{BaCl2} 溶液。（\ce{CaCO3} 溶于盐酸放出 \ce{CO2}；若沉淀中有 \ce{CaSO4}，溶解后加 \ce{BaCl2} 会生成 \ce{BaSO4} 白色沉淀，无沉淀说明无 \ce{CaSO4}。）(2) $n(\ce{CaSO4})=1.00/136\approx7.4\times10^{-3}$ mol，$n(\ce{CO3^2-})=0.100\times0.100=1.00\times10^{-2}$ mol，\ce{CaSO4} 不足；反应Ⅰ的平衡常数很大，\ce{CaSO4} 可完全转化。pH 不变只是因为一种反应物耗尽，并非平衡。(3) 要保证 \ce{CaSO4} 有剩余：按完全反应计算至少需 $0.0100\times136=1.36$ g，故将不少于 1.36 g 硫酸钙加入 100.0 mL 0.100 mol·L$^{-1}$ \ce{Na2CO3} 溶液中，在 25 ℃、搅拌条件下用 pH 计测 pH；当 pH 不变时，加入 \ce{Na2SO4} 固体，若 pH 增大（\ce{SO4^2-} 浓度增大使反应Ⅰ逆移，$c(\ce{CO3^2-})$ 增大，碱性增强），说明反应Ⅰ的平衡已建立（体现了“平衡可以被打破并移动”这一特征）。
\item (1) 设 \ce{N2} 转化 $x$ mol：平衡时 \ce{N2} $(1-x)$、\ce{H2} $(3-3x)$、\ce{NH3} $2x$，总 $(4-2x)$ mol。$\dfrac{2x}{4-2x}=0.25$，$x=0.4$；$\alpha(\ce{H2})=\dfrac{1.2}{3}=40\%$。(2) $n_\text{总}=3.2$ mol，物质的量分数：\ce{N2} 0.1875、\ce{H2} 0.5625、\ce{NH3} 0.25。$K_p=\dfrac{(0.25\times20)^2}{(0.1875\times20)\times(0.5625\times20)^3}\approx4.68\times10^{-3}$ MPa$^{-2}$。(3) 小于。恒容时随反应进行气体分子数减少、压强降低，相当于在比恒压时更低的压强下建立平衡，平衡转化率较低。
\item (1) $m_1>m_2>m_3$：\ce{H2} 相对越多，\ce{CO2} 的转化率越高。(2) 正反应放热，升温平衡逆向移动，\ce{CO2} 平衡转化率下降。(3) $K_a>K_b$（放热反应，温度越高 $K$ 越小；$K$ 只与温度有关，与投料比无关）。
\item (1) 设乙苯 1 mol，平衡时乙苯 0.6 mol、苯乙烯 0.4 mol、\ce{H2} 0.4 mol，$n_\text{总}=1.4$ mol。$K_p=\dfrac{(\frac{0.4}{1.4}p)^2}{\frac{0.6}{1.4}p}=\dfrac{0.16}{0.84}\times100\approx19$ kPa。(2) 设转化率为 $\alpha$，$n_\text{总}=10+\alpha$：$K_p=\dfrac{\alpha^2p}{(1-\alpha)(10+\alpha)}=19.05$，整理得 $119.05\alpha^2+171.45\alpha-190.5=0$，$\alpha\approx0.74$，即约 74\%（远高于 40\%）。(3) ①恒压下水蒸气稀释，降低各反应物、产物的分压，相当于减压，平衡向气体分子数增大的正方向移动，提高乙苯转化率；②水蒸气热容大，可为吸热反应供热、维持反应温度；③水蒸气可与催化剂表面积炭反应（\ce{C + H2O -> CO + H2}），减少积炭，延长催化剂寿命。
\item ①$(1-y_2)\times100\%$。起始 $n_2/n_0=1$，$t_0$ 时为 $y_2$，转化了 $n_0(1-y_2)$。②设反应(e)消耗 \ce{N2H4} $n_\text{e}$，反应(f)消耗 $n_\text{f}$：$n_\text{e}+n_\text{f}=n_0(1-y_2)$；(e)生成 $n_\text{e}$ mol \ce{N2} 和 $2n_\text{e}$ mol \ce{H2}，(f)生成 $\frac13n_\text{f}$ mol \ce{N2}：$3n_\text{e}+\frac13n_\text{f}=y_1n_0$。将第一式乘 $\frac13$ 后与第二式相减：$\frac83n_\text{e}=n_0\left(y_1-\frac{1-y_2}{3}\right)$，$n_\text{e}=\dfrac{(3y_1+y_2-1)n_0}{8}$。选择性 $=\dfrac{n_\text{e}}{n_0(1-y_2)}\times100\%=\dfrac{3y_1+y_2-1}{8-8y_2}\times100\%$。
\item ①0.82。303 K、100 kPa 饱和时每摩尔 D 吸收 $\alpha+\beta=0.15+0.86=1.01$ mol；350 K、40 kPa 平衡时为 $0+0.19=0.19$ mol；释放 $1.01-0.19=0.82$ mol。②0.22。350 K 时 $\alpha=0$，HCl 只以 \ce{D*$x$HCl} 形式被吸收。$p=100$ kPa 时 $\beta=0.25$：$n(\mathrm{DH})=\dfrac{0.25}{x}$ mol，$n(\mathrm{D})=1-\dfrac{0.25}{x}$，$K=\dfrac{0.25/x}{(1-0.25/x)\times1^x}=1.0$，解得 $\dfrac{0.25}{x}=0.5$，$x=0.5$。$p=64$ kPa 时：$\dfrac{n(\mathrm{DH})}{[1-n(\mathrm{DH})]\times0.64^{0.5}}=1.0$，即 $n(\mathrm{DH})=0.8[1-n(\mathrm{DH})]$，$n(\mathrm{DH})=\dfrac49$ mol；吸收 HCl 的量 $=\beta=x\cdot n(\mathrm{DH})=0.5\times\dfrac49\approx0.22$ mol。
\item (1) 设生成 \ce{N2O4} 的分压为 $x$：$p(\ce{NO2})=100-2x$，总压 $100-x=70$，$x=30$ kPa；$p(\ce{NO2})=40$ kPa，$p(\ce{N2O4})=30$ kPa，$K_p=\dfrac{30}{40^2}=0.01875$ kPa$^{-1}$。(2) 平衡时 $\vz=\vn$，$K_p=\dfrac{\kz}{\kn}$。(3) 此时 $p(\ce{N2O4})=\dfrac{100-60}{2}=20$ kPa，$\dfrac{\vz}{\vn}=\dfrac{\kz\times60^2}{\kn\times20}=0.01875\times180=3.375$（$>1$，反应正向进行）。(4) 升温后 \ce{NO2} 的体积分数增大，说明平衡逆向移动，逆反应吸热，正反应放热，$\Delta H<0$。
\item (1) $<$；低温、高压。升温 $x$ 减小，平衡逆移，正反应放热；加压 $x$ 增大，有利于气体分子数减小的正反应。(2) 设起始 CO 为 1 mol，生成 \ce{Ni(CO)4} $y$ mol：剩余 CO $(1-4y)$ mol，气体总量 $(1-3y)$ mol。$\dfrac{y}{1-3y}=0.90$，$y=\dfrac{0.90}{3.7}\approx0.243$；$\alpha(\ce{CO})=4y\approx97.3\%$。$p(\ce{Ni(CO)4})=0.90$ MPa，$p(\ce{CO})=0.10$ MPa，$K_p=\dfrac{0.90}{0.10^4}=9.0\times10^3$ MPa$^{-3}$。(3) 低温时平衡正移，粗镍中的 Ni 转化为气态 \ce{Ni(CO)4} 与杂质分离；高温时平衡逆移，\ce{Ni(CO)4} 分解析出纯镍，CO 可循环使用——利用温度对平衡的影响实现“搬运”提纯。
\end{jiexilist}
